\input{../preamble.tex}

\SetDate[17/04/2026]

\usepackage{minted}

\def\arraystretch{2}
\setlength\tabcolsep{18pt}

\begin{document}
\pagestyle{fancy}
\fancyhead[L]{Seconde 5}
\fancyhead[C]{\textbf{Méthode de dichotomie}}
\fancyhead[R]{\today}

\section*{Le Juste Prix}

\begin{namedtheorem}[Règles du jeu]
	Le joueur doit deviner le prix d'un objet pour le gagner.
	Pour cela, il peut faire plusieurs estimations.
	
	Le présentateur connaît le prix de l'objet. 
	À chaque estimation, il répond « C'est plus ! » ou \mbox{« C'est moins ! »} selon que le prix de l'objet est plus grand ou plus petit que l'estimation donnée.
	
	Le joueur est chronométré et doit trouver le prix exact de l'objet le plus vite possible.
\end{namedtheorem}

%\begin{namedtheorem}[Simplification]
	On suppose que le prix est un nombre entier inférieur ou égal à 1000€.
%\end{namedtheorem}

\subsection*{Premier objet}

\hspace{-40pt}
\begin{tabular}{|*{11}{c|}}\hline
	Essai n° & 1 & 2 & 3 & 4 & 5 & 6 & 7 & 8 & 9 & 10 
	\\\hline
	Estimation & & & & & & & & & &
	\\\hline
	Comparaison & & & & & & & & & &
	\\\hline
\end{tabular}

\vspace{10pt}
Prix final trouvé :
\hspace{6cm}
Nombre d'estimations :

\subsection*{Deuxième objet}

\hspace{-40pt}
\begin{tabular}{|*{11}{c|}}\hline
	Essai n° & 1 & 2 & 3 & 4 & 5 & 6 & 7 & 8 & 9 & 10 
	\\\hline
	Estimation & & & & & & & & & &
	\\\hline
	Comparaison & & & & & & & & & &
	\\\hline
\end{tabular}

\vspace{10pt}
Prix final trouvé :
\hspace{6cm}
Nombre d'estimations :

%\begin{namedtheorem}[Oracle]
%	Vincent Lagaf.
%	
%	Il connait le prix exact et nous dit s'il est plus grand ou plus petit que le prix proposé.
%\end{namedtheorem}
%
%\begin{namedtheorem}[Formalisation]
%	On cherche un prix $P \in [a ; b]$.
%	On ne peut que demander si un prix $x$ vérifie $x \leq P$ ou $x \geq P$.
%\end{namedtheorem}
%
%\begin{namedtheorem}[Stratégie optimale]
%	À chaque étape, on divise l'intervalle en deux et on appelle l'oracle.
%	Ceci permet d'exclure la moitié des valeurs.
%	
%	La stratégie est optimale car si l'oracle agit comme adversaire, il peut choisir de changer $P$ selon le prix $x$ proposé par le joueur.
%\end{namedtheorem}

\section*{À la recherche de racine}

\begin{namedtheorem}[Nomenclature]
	On appelle \emph{racine} de $f$ un antécédent $x$ qui vérifie $f(x) = 0$.
\end{namedtheorem}

\begin{namedtheorem}[But]
	Trouver une racine de $f$ : un $x$ du domaine tel que $f(x) = 0$.
\end{namedtheorem}

\begin{namedtheorem}[Hypothèses]
	Le tracé de $\C_f$ est continu : on la trace sans lever son stylo.
	
	Sans perte de généralité, $f(x)$ passe du négatif au positif autour de sa racine (on peut prendre $-f$ dans le cas contraire).
\end{namedtheorem}

\begin{namedtheorem}[Oracle]
	Étant donné un $x$, un appel à l'oracle permet de connaître le signe de $f(x)$.
\end{namedtheorem}

\begin{namedtheorem}[Stratégie optimale]
	$f(x) = x^2 - 7$
	$a = 2$
	$b = 3$
	$\text{precision} = 10^{-5}$
	
	Tant que \\
	Faire : \\
		Calcul du milieu $x = \frac{a+b}2$ \\
		Appel à l'oracle pour connaitre le signe de $f(x)$ \\
		Si $f(x) > 0$ \\
		$b = x$ \\
		Sinon, \\
		$a = x$ \\
\end{namedtheorem}

%Différence avec le balayage.

%\subsubsection*{Exemple 1}
%\python{while}
% Commentaires :
% 1. On initialise x = 0
% 2. Tant que x est strictement inférieur à 2
% 3. Remplacer x par sa valeur + 0,1
% 4. À la fin de la boucle, imprimer la valeur de x
% valeur attendue : 2
%
%Impression finale :
%
%Nombre d'itérations de la boucle \texttt{while} :
%
%\subsubsection*{Exemple 2}
%\python{while2}
% Commentaires :
% 1. On initialise x = 0
% 2. Tant que x est positif ou nul à 2
% 3. Remplacer x par sa valeur + 0,1
% 4. À la fin de la boucle, imprimer la valeur de x
% valeur attendue : boucle infinie !
%
%Impression finale :
%
%Nombre d'itérations de la boucle \texttt{while} :
%
%\newpage
%
%\exe{,difficulty=1}{
%	Compléter et implémenter le programme de balayage ci-dessous.
%	Que renvoie-t-il ?
%}{exe:balayage4}{
%	%\python{balayage-solution}
%}
%
%\python{balayage}
%
%
%\exe{}{
%	Donner la valeur de $\sqrt7$ à $10^{-5}$ près à l'aide d'un programme de balayage.
%}{exe:precision}{
%	En remplaçant la variable \texttt{amplitude} par \texttt{0.00001} ou \texttt{1e-5}, on obtient
%		\[ 2,64575 < \sqrt7 < 2,64576. \]
%	%\python{balayage-solution2}
%}
%
%\exe{}{
%	Donner la valeur de $\sqrt{700}$ à $10^{-4}$ près
%		\begin{enumerate}[label=\roman*)]
%			\item
%			à l'aide d'un programme de balayage ; et
%			\item
%			en montrant que $\sqrt{700} = 10\sqrt7$ et en utilisant l'exercice \ref{exe:precision}.
%		\end{enumerate}
%}{exe:precision2}{
%	En remplaçant la fonction $f$ par $f(x) = x^2 - 700$, et l'amplitude par $10^{-4}$, on obtient
%		\[ 26,4575 < \sqrt{700} < 26,4576. \]
%	%\python{balayage-solution3}
%	Comme $\sqrt{700} = \sqrt{100\times7} = \sqrt{100}\times\sqrt{7}=10\sqrt7$, on peut multiplier l'encadrement de l'exercice \ref{exe:precision} par 10 pour conclure immédiatement.
%}

\subsection*{Trouver un polynôme annulateur}


La fonction $f(x) = x^2 - 3$ 
	\begin{enumerate}
		\item est un \emph{polynôme}, car elle s'exprime comme somme de multiples de $1, x, x^2, x^3, \hdots$ ; 
		\item est \emph{annulatrice} de $\sqrt3$ car $f\bigl(\sqrt3\bigr) = 0$ ;
		\item est à \emph{coefficients entiers} car les multiples de $1, x, x^2$ sont $-3 ; 0 ; 1$ respectivement et sont entiers.
	\end{enumerate}
	
\exe{, difficulty=0}{
	Trouver un polynôme à coefficients entiers et annulateur de $\sqrt2$.
}{exe:p-annul-sqrt2}{
	$f(x) = x^2 -2$ fonctionne, mais pas $f(x) = x-\sqrt2$ !
}
	
\exe{, difficulty=1}{
	Trouver un polynôme à coefficients entiers et annulateur de $\dfrac{\sqrt3}2$.
}{exe:p-annul-sqrt32}{
	$g(x) = 4x^2 -3$ fonctionne, car $\frac{\sqrt3}2 = \sqrt{\frac34}$.
}
	
\exe{, difficulty=2}{
	Trouver un polynôme à coefficients entiers et annulateur à la fois de $\sqrt2$ et de $\dfrac{\sqrt3}2$.
}{exe:p-annul-sqrt2-sqrt32}{
	$h(x) = f(x)\cdot g(x) = (x^2 - 2)(4x^2 -3) = 4x^4 - 11x^2 + 6$ fonctionne, avec $f$ et $g$ des deux exercices précédents.
}


\exe{, difficulty=3}{
	Trouver un polynôme à coefficients entiers et annulateur de
		$ \sin(36\degree) = \sqrt{\dfrac{5-\sqrt5}8}. $

}{exe:p-annul}{
	En nommant $x = \sqrt{\dfrac{5-\sqrt5}8}$, la mise au carré donne
		\[ x^2 =  \dfrac{5-\sqrt5}8. \]
	En multipliant par 8 et en soustrayant 5, on trouve
		\begin{align*}
			 8x^2 &= 5 - \sqrt5, \\
			 8x^2 - 5 &= - \sqrt5.
		\end{align*}
	Une nouvelle mise au carré nous permet d'obtenir les égalités suivantes.
		\begin{align*}
			(8x^2 - 5)^2 &= 5 \\
			64x^4 - 80x^2 + 25 &= 5 \\
			64x^4 - 80x^2 + 20 &= 0 \\
			16x^4 - 20x^2 + 5 &= 0
		\end{align*}
	Un polynôme annulateur est donc $p(x) = 16x^4 - 20x^2 + 5$.
}

\begin{namedtheorem}[Théorème]
	$\pi$ est un nombre transcendant : il n'existe pas de polynôme à coefficients entiers annulateur de $\pi$.
\end{namedtheorem}

\subsection*{Étude de complexité}


\exe{, difficulty=1}{
	Montrer que l'algorithme par dichotomie avec fonction $f(x) = x^2 - 2$ et domaine de départ $[1 ; 2]$ n'atteint jamais la racine.
	
	\emph{Indication : montrer que les bornes stockées sont toujours rationnelles.}
}{exe:appf1}{
}

\exe{, difficulty=2}{
	En partant d'un encadrement à l'unité d'une racine, combien d'étapes faut-il pour calculer ses 10 premières décimales ?
	
	Montrer qu'on cherche le plus petit $n\in\N$ tel que $2^{n} \geq 10^{10}$
	puis résoudre ce problème algorithmiquement et vérifier qu'on obtient bien $n=34$.
}{exe:appf4}{
}


%%%%%%%%%%%

\newpage
\fancyhead[C]{\textbf{Solutions}}
\shipoutAnswer

\end{document}
